\chapter*{Abstract}

This project is part of the development of MICROLAB, a virtual laboratory for 
teaching embedded systems that uses QEMU as an emulation platform to replicate 
boards based on STM32F4 microcontrollers used at the UPM (Technical University of 
Madrid). The main objective was to incorporate the I$^2$C bus into the STM32F405 
microcontroller and emulate one of the sensors used in laboratory exercises, the 
LIS3DH triaxial accelerometer, so that the firmware developed with STM32CubeIDE 
could run without modification on a completely virtual platform. To achieve this, 
the architectures of the STM32F405 and the LIS3DH, their register maps, and the 
interaction at the I$^2$C protocol level were analyzed in detail, defining a design 
that prioritizes register-level fidelity and compatibility with ST's HAL library.

On the microcontroller side, a complete model of the STM32F405 I$^2$C controller 
has been implemented in QEMU. This model is based on a 
finite state machine that replicates the START, address, read/write, and STOP 
sequences, as well as the generation of event and error interrupts (EV/ER) and the 
management of the CR1, CR2, SR1, SR2, and DR registers. On the sensor side, a 
detailed model of the LIS3DH as an I$^2$C slave device has been 
developed, implementing its 64-register map, the three operating modes (normal, high 
resolution, and low power), the configurable measurement ranges ($\pm$2~g, $\pm$4~g, 
$\pm$8~g, $\pm$16~g), and the generation of the data-ready interrupt (DRDY) on the 
INT1 line. The model also exposes QOM properties (\texttt{accel-x/y/z}, 
\texttt{temp}) that allow for the controlled injection of acceleration and 
temperature data during testing, facilitating the systematic verification of the 
emulated device's behavior.

Both models have been integrated into a QEMU machine based on the Netduino Plus 2 
board, which incorporates the STM32F405 SoC, peripheral address mapping, and correct 
routing of the NVIC and EXTI interrupt lines, connecting the LIS3DH as a slave on 
the I$^2$C1 bus at address \texttt{0x18}. Real firmwares developed with 
STM32CubeIDE and the HAL library were run on this platform, verifying I$^2$C bus 
initialization, accelerometer register configuration, auto-incrementing multi-byte 
readings, and the reception of data-ready interrupts. This demonstrates functional 
compatibility with real hardware in typical teaching scenarios. Overall, this work 
provides MICROLAB with a complete emulation of the STM32F405--LIS3DH I$^2$C 
subsystem, enabling new virtual sensing and communication exercises without the need 
for physical hardware and establishing a reusable foundation for the future 
integration of other microcontrollers and peripherals from the STM32F4 family.


\section*{Keywords}
QEMU, HAL, STM32, ARM, LIS3DH, emulation, virtualization, I$^2$C